% 02_VRL_CoJeVyrok.tex
% Rozpracovana priprava: 1. rocnik, hodina 02, Vyrokova logika.
% Kompilace z Work/MAT: latexmk -cd -pdf 1_rocnik/pripravy/02_VRL_CoJeVyrok.tex
\documentclass[11pt,a4paper]{article}

% Cesta ke stylu je relativni k souboru (latexmk s prepinacem -cd).
% Spolecny styl pro ucitelske pripravy MAT (Work/MAT)
% Verze 1.1, 2026-10-23
\NeedsTeXFormat{LaTeX2e}
\ProvidesPackage{priprava}[2026/10/23 v1.1 Sablona ucitelskych priprav MAT]

% Zakladni typografie a cestina
\RequirePackage[T1]{fontenc}
\RequirePackage[utf8]{inputenc}
\RequirePackage[czech]{babel}
\RequirePackage{lmodern}
\RequirePackage{microtype}
\RequirePackage{amsmath,amssymb}
\RequirePackage{xcolor}
\RequirePackage{enumitem}
\RequirePackage[most]{tcolorbox}
\RequirePackage{tabularx,booktabs,array}
\RequirePackage{fancyhdr}
\RequirePackage{needspace}
\RequirePackage{geometry}
\RequirePackage[hidelinks]{hyperref}

\geometry{a4paper,top=17mm,bottom=18mm,left=23mm,right=23mm}
\definecolor{notegray}{gray}{0.39}
\setlength{\parindent}{0pt}
\setlength{\parskip}{2.5pt plus 1pt minus 0.5pt}
\setlist[itemize]{topsep=2pt,itemsep=1pt,parsep=0pt,leftmargin=1.45em}
\setlist[enumerate]{topsep=2pt,itemsep=2pt,parsep=0pt,leftmargin=1.65em}
\setcounter{secnumdepth}{0}
\renewcommand{\arraystretch}{1.12}

% Zahlavi: vlevo rocnik | globalni tema, vpravo cislo: nazev hodiny
% Zapati zustava zcela prazdne.
\pagestyle{fancy}
\fancyhf{}
\renewcommand{\headrulewidth}{0.35pt}
\renewcommand{\footrulewidth}{0pt}
\setlength{\headheight}{14pt}
\newcommand{\nastavPripravu}[4]{%
  \fancyhead[L]{\small #1\enspace\textbar\enspace #2}%
  \fancyhead[R]{\small #3: #4}%
}

% Minimalisticke cernobile boxy. Kazda definice zvlast.
\tcbset{commonlesson/.style={enhanced,breakable=false,colback=white,colframe=black,
  colbacktitle=white,coltitle=black,boxrule=0.45pt,
  arc=0.8mm,outer arc=0.8mm,left=3mm,right=3mm,top=1.6mm,bottom=1.6mm,
  toptitle=1mm,bottomtitle=0.7mm,fonttitle=\bfseries,
  attach boxed title to top left={xshift=2.5mm,yshift=-1.15mm},
  boxed title style={colback=white,colframe=white,boxrule=0pt,left=1mm,right=1mm}}}
\newtcolorbox{zadaniBox}[1][]{commonlesson,title={Zadání},#1}
\newtcolorbox{prikladBox}[1][]{commonlesson,title={Příklad},#1}
\newtcolorbox{definiceBox}[1][]{commonlesson,title={Definice},#1}

% Spolecny box pro zadani a jeho reseni, bez mezery mezi castmi.
% Pouziti:
%   \begin{prikladReseniBox}
%     Zadani ulohy.
%     \tcblower
%     \begin{align*} ... \end{align*}
%   \end{prikladReseniBox}
% Spodni cast ma jemny sedy podklad a delici linku; lze v ni sazet
% vice radku matematiky. Obe casti maji jen jeden spolecny ramecek.
\tcbset{solutionlesson/.style={
  commonlesson,
  enhanced jigsaw,
  breakable,
  colbacklower=black!3,
  segmentation style={solid,draw=black!25,line width=0.35pt},
  fontlower=\small\color{notegray},
  before lower={\noindent\textbf{Řešení:}\enspace},
  leftlower=3mm,rightlower=3mm
}}
\newtcolorbox{prikladReseniBox}[1][]{solutionlesson,title={Příklad},#1}
\newtcolorbox{zadaniReseniBox}[1][]{solutionlesson,title={Zadání},#1}

% Pomucky pro autora pripravy: sedive poznamky; \reseni pro starsi pripravy
\newcommand{\poznamka}[1]{\par\begingroup\small\color{notegray}\itshape
  \textnormal{\textbf{Poznámka pro učitele.}}\enspace #1\par\endgroup}
\newcommand{\reseni}[1]{\par\begingroup\small\color{notegray}
  \textbf{Očekávané řešení:}\enspace #1\par\endgroup}
\newcommand{\blok}[1]{\needspace{5\baselineskip}\subsection*{#1}}
\newcommand{\vyrok}[1]{\uv{#1}}

\nastavPripravu{1. ročník}{Výroková logika}{02}{Co je výrok?}

\begin{document}

\blok{1. Motivace}
Na tabuli napiš následující čtyři věty a nech žáky hlasovat, zda jim lze přiřadit pravdivostní hodnotu.
\begin{prikladReseniBox}
\begin{enumerate}[label=\alph*)]
  \item Sedm je prvočíslo.
  \item Otevřete si sešity.
  \item $3+4=9$.
  \item Matematika je nejzábavnější předmět.
\end{enumerate}
\tcblower
a) pravdivý výrok; b) není výrok (rozkaz); c) nepravdivý výrok; d) bez upřesněného objektivního kritéria nejde o jednoznačně pravdivé či nepravdivé sdělení.
\end{prikladReseniBox}
\poznamka{Nezaměňovat \uv{nepravdivý výrok} a \uv{není výrok}. Pro žáky je to zásadní rozdíl.}

\blok{2. Co je výrok}
\begin{definiceBox}
\textbf{Výrok} je sdělení, u něhož má smysl rozhodovat, zda je \textbf{pravdivé}, nebo \textbf{nepravdivé}. Jeho pravdivostní hodnotu značíme zpravidla $1$ (pravda) a $0$ (nepravda).
\end{definiceBox}

Výroky můžeme označovat malými písmeny, například
\[
  p:\; 12\text{ je dělitelné }3,\qquad q:\; 5<2.
\]
Výrok $p$ má hodnotu $1$, výrok $q$ má hodnotu $0$.

\poznamka{Pojem \uv{výrok} se vztahuje k obsahu sdělení, nikoliv pouze k tomu, zda má věta tečku. Také matematický zápis může být výrokem.}

\blok{3. Rozlišování výroků}
\begin{zadaniReseniBox}
Rozhodni, které položky jsou výroky. U výroků urči pravdivostní hodnotu.
\begin{enumerate}[label=\alph*)]
  \item $2^4=16$.
  \item Kolik je odmocnina ze čtyř?
  \item Každé sudé celé číslo je dělitelné dvěma.
  \item $x+5=8$, kde $x$ není určeno.
  \item Existuje prvočíslo větší než 100.
  \item Ať už je přestávka!
\end{enumerate}
\tcblower
a) výrok, $1$; b) otázka; c) výrok, $1$; d) otevřená formule, bez určení $x$ není výrok; e) výrok, $1$; f) přání.
\end{zadaniReseniBox}

\blok{4. Atomický, složený výrok a formule}
\begin{definiceBox}
\textbf{Atomický výrok} je výrok, který nelze rozložit na jednodušší výroky pomocí výrokových spojek.
\end{definiceBox}

\begin{definiceBox}
\textbf{Složený výrok} vzniká užitím výrokových spojek na jeden či více výroků (např. \uv{a}, \uv{nebo}, \uv{jestliže -- pak}).
\end{definiceBox}

\begin{prikladBox}
$p$: \vyrok{4 je sudé číslo.}\qquad $q$: \vyrok{9 je dělitelné třemi.}
\begin{itemize}
  \item $p$ je atomický výrok.
  \item \vyrok{4 je sudé číslo a 9 je dělitelné třemi} je složený výrok $p\land q$.
\end{itemize}
\end{prikladBox}

\begin{zadaniReseniBox}
U každé položky řekni, zda jde o atomický výrok, složený výrok, nebo o výraz, který zatím není výrokem.
\begin{enumerate}[label=\alph*)]
  \item $15$ je liché číslo.
  \item $15$ je liché číslo a $15$ je dělitelné pěti.
  \item $a^2+b^2=c^2$ pro neurčená $a,b,c$.
  \item Jestliže je číslo dělitelné deseti, pak je dělitelné pěti.
\end{enumerate}
\tcblower
a) atomický; b) složený; c) otevřená formule s neurčenými proměnnými; d) složený výrok.
\end{zadaniReseniBox}
\poznamka{\uv{Formule} není synonymum pro \uv{složený výrok}: otevřená formule s volnými proměnnými může teprve po dosazení nebo kvantifikaci vyjadřovat výrok. Nezabíhat do formální syntaxe.}

\blok{5. Výrokové spojky -- zápis a čtení}
\begin{center}
\small
\begin{tabularx}{\linewidth}{@{}p{36mm}p{23mm}X@{}}
\toprule
Název & Zápis & Jak čteme\\
\midrule
Negace & $\neg p$ & není pravda, že $p$\\
Konjunkce & $p\land q$ & $p$ a zároveň $q$\\
Disjunkce & $p\lor q$ & $p$ nebo $q$ (případně obojí)\\
Implikace & $p\Rightarrow q$ & jestliže $p$, pak $q$\\
Ekvivalence & $p\Leftrightarrow q$ & $p$ právě tehdy, když $q$\\
\bottomrule
\end{tabularx}
\end{center}

\begin{zadaniReseniBox}
Nechť $p$: \vyrok{Prší.} a $q$: \vyrok{Mám deštník.} Přečti nahlas výroky
\[
  p\land q,\qquad p\lor q,\qquad p\Rightarrow q,\qquad p\Leftrightarrow q,\qquad \neg p.
\]
\tcblower
Například $p\Rightarrow q$: \uv{Jestliže prší, pak mám deštník.} Ostatní čteme analogicky podle tabulky.
\end{zadaniReseniBox}
\poznamka{Upozorni, že matematické \uv{nebo} je zpravidla nevýlučné. Zatím po žácích nechci všechny pravdivostní hodnoty složených výroků ani jejich tabulky.}

\blok{6. Závěr a rychlé ověření}
\begin{zadaniReseniBox}
\begin{enumerate}
  \item Je \vyrok{$x>0$} bez dalšího určení výrokem? Proč?
  \item Zapiš symbolem: \uv{Sedm je prvočíslo a osm je sudé číslo.}
  \item Jak se čte $p\Leftrightarrow q$?
\end{enumerate}
\tcblower
1. Ne, pravdivost závisí na $x$. 2. $p\land q$, kde $p$: 7 je prvočíslo, $q$: 8 je sudé. 3. $p$ právě tehdy, když $q$.
\end{zadaniReseniBox}

\textbf{Shrnutí na tabuli:} \;Výrok $\to$ pravdivostní hodnota $\to$ atomický/složený výrok $\to$ logické spojky.

\textbf{Domácí procvičení (volitelné):} Napiš tři výroky, dva nevýroky a dva složené výroky. U složených výroků podtrhni spojku.

\end{document}
